\documentclass[11pt,a4paper]{article}
\usepackage[margin=26mm,headheight=15pt]{geometry}
\usepackage[T1]{fontenc}
\usepackage[utf8]{inputenc}
\usepackage{libertinus}
\usepackage{amsmath,amssymb,amsthm,mathtools}
\usepackage{microtype,booktabs,array,longtable}
\usepackage{xcolor,enumitem,fancyhdr,titlesec,xurl}
\usepackage{aliascnt}
\usepackage{hyperref}
\usepackage[nameinlink,noabbrev]{cleveref}
\definecolor{ink}{RGB}{25,53,75}
\definecolor{muted}{RGB}{85,90,96}
\hypersetup{colorlinks=true,linkcolor=ink,citecolor=ink,urlcolor=ink,
 pdftitle={Beyond the Square-Root Cusp},
 pdfauthor={Research draft prepared with ChatGPT for Vladimir Reshetnikov},
 pdfsubject={Logarithmic critical corrections and equilibrium selection for digital-product pressure}}
\titleformat{\section}{\Large\bfseries\color{ink}}{\thesection}{.7em}{}
\titleformat{\subsection}{\large\bfseries}{\thesubsection}{.7em}{}
\pagestyle{fancy}\fancyhf{}
\fancyhead[L]{\small Beyond the square-root cusp}
\fancyhead[R]{\small ProveIt research}
\fancyfoot[C]{\thepage}
\setlength{\emergencystretch}{3em}
\setlength{\parskip}{.2em}
\setlist{itemsep=.25em,topsep=.5em}
\allowdisplaybreaks
\numberwithin{equation}{section}
\newtheorem{theorem}{Theorem}[section]
\newaliascnt{lemma}{theorem}\newtheorem{lemma}[lemma]{Lemma}\aliascntresetthe{lemma}
\newaliascnt{proposition}{theorem}\newtheorem{proposition}[proposition]{Proposition}\aliascntresetthe{proposition}
\newaliascnt{corollary}{theorem}\newtheorem{corollary}[corollary]{Corollary}\aliascntresetthe{corollary}
\theoremstyle{definition}
\newaliascnt{definition}{theorem}\newtheorem{definition}[definition]{Definition}\aliascntresetthe{definition}
\newtheorem{question}{Research question}
\theoremstyle{remark}
\newaliascnt{remark}{theorem}\newtheorem{remark}[remark]{Remark}\aliascntresetthe{remark}
% ed. (2026-09-29): unnumbered environment for editorial notes.
\newtheorem*{ednote}{Editorial note (ProveIt, 2026-09-29)}
\newcommand{\TT}{\mathbb T}
\newcommand{\RR}{\mathbb R}
\newcommand{\CC}{\mathbb C}
\newcommand{\NN}{\mathbb N}
\newcommand{\ZZ}{\mathbb Z}
\newcommand{\dd}{\,\mathrm d}
\newcommand{\Lop}{\mathcal L}
\newcommand{\Id}{\mathrm I}
\newcommand{\Proj}{\Pi}
\newcommand{\spec}{\operatorname{spec}}
\newcommand{\dist}{\operatorname{dist}}
\newcommand{\norm}[1]{\left\lVert#1\right\rVert}
\newcommand{\abs}[1]{\left|#1\right|}
\newcommand{\code}[1]{\texttt{\detokenize{#1}}}
\newcommand{\Holder}{C^\alpha(\TT)}
\newcommand{\Leb}{\mathfrak m}
\newcommand{\bconst}{B_b}
\newcommand{\rhoess}{r_{\mathrm{ess}}}

\begin{document}
\hypersetup{pageanchor=false}
\begin{titlepage}
\thispagestyle{empty}
\setlength{\parindent}{0pt}
{\small\sffamily\color{ink} PROVEIT\quad /\quad ANALYSIS AND DYNAMICAL SYSTEMS}\par
\vspace{15mm}
{\Huge\bfseries\color{ink} Beyond the\\[3mm] Square-Root Cusp\par}
\vspace{7mm}
{\Large Logarithmic critical corrections and equilibrium\\[2mm]
selection for digital-product pressure\par}
\vspace{9mm}
{\large Research draft prepared for Vladimir Reshetnikov\par}
{Prepared with ChatGPT\quad\textbullet\quad 29 September 2026\par}
\vspace{7mm}\hrule\vspace{3mm}
\begin{abstract}
The ProveIt critical-pressure article determines a square-root singularity
and explicitly asks for its next term and for the limiting equilibrium
measures. We address both questions. For the consecutive digital mask in
any integer base $b\ge2$, put $\kappa_b=\sqrt{2(b-1)\log b}$. At the
critical absolute-moment exponent, the unnormalized cylinder pressure is
\[
\begin{split}
 P_{b,1}(c)={}&\kappa_b\sqrt{|c|}-(b-1)|c|\log(1/|c|)\\
 &+\bigl[b\log b+(b-1)(\gamma-1)\bigr]|c|
 +O_{b,\eta}(|c|^{3/2-\eta})
\end{split}
\]
for every $0<\eta<1/2$. Thus both a logarithmic term and an ordinary
linear term occur. The logarithm comes from an explicitly evaluated
finite-part response functional, rather than from a formal continuation
of a supercritical zeta coefficient.

At nonzero phases near zero there is a unique equilibrium probability.
At the critical exponent it converges weakly to
$\tfrac12\delta_0+\tfrac12\Leb$, where $\Leb$ is normalized Lebesgue
measure. More generally, along $s=1+u\sqrt{|c|}$, bounded detuning selects
an explicit mixture whose atomic mass is
\[
 \frac12\left(1+
 \frac{u\log b}{\sqrt{u^2(\log b)^2+8(b-1)\log b}}\right).
\]
We obtain the individual left and right eigenvector asymptotics, a
$\sqrt{|c|}$-scale boundary layer, and the exact leading divergence of a
finite-product amplitude. The spectral foundation and the leading
crossover are credited to the repository predecessor and reconstructed
as needed. Detailed proofs, reproducible diagnostics, a claims audit,
and further research questions accompany the article.
\end{abstract}
\vfill
{\small\color{muted} Status: ordinary mathematical proofs in an unrefereed research draft;
not Lean/Rocq verified. Numerical data are not interval certificates.
No unrestricted claim of worldwide priority is made.}
\end{titlepage}
\hypersetup{pageanchor=true}
\begingroup\small\setlength{\parskip}{0pt}\tableofcontents\endgroup
\newpage

\section{The precise research target}
\label{sec:intro}

\subsection{What was already known in the inspected repository}
The immediate predecessor is \emph{Critical Cusps in Digital-Product
Pressure} \cite{CriticalRepo}. It proves the leading law
$P_{b,1}(c)\sim\kappa_b\sqrt{|c|}$, constructs a size-two Jordan block at
the atomic phase, and obtains the leading bounded-detuning and finite-size
crossover. Those results are \emph{not} claimed anew here. Its further
research agenda contains two particularly concrete questions:
\begin{enumerate}[label=\arabic*.,leftmargin=2em]
\item After the square root, does the pressure have a
$|c|\log|c|$ term, a $|c|$ term, or both, and what are their coefficients?
\item Which normalized eigenvectors and equilibrium measures emerge as
the phase tends to the atomic phase?
\end{enumerate}
These are the predecessor's questions entitled ``The next critical term''
and ``Eigenfunctions, eigenmeasures, and equilibrium measures.'' The
present article answers the first with an explicit remainder and answers
the weak-selection and first boundary-layer parts of the second.

The earlier \emph{Fractional Cusps at the Atomic Phase} treats exponents
$s>1$ \cite{FractionalRepo}. A later subcritical article treats all
$0<s<1$ \cite{SubcriticalRepo}. Neither the leading square root nor the
subcritical linear response is a new contribution of this package.
This distinction matters in a rapidly expanding research corpus.

The primary literature studies generalized Thue--Morse measures and
phase-dependent singular potentials \cite{GKS,GLS}. In particular,
\cite[Theorems 2.2--2.3 and the following discussion]{GLS} proves
phase continuity at nonnegative pressure orders and asks about stronger
regularity at general positive orders. Our result is a refined local
answer at the exceptional atomic phase; it is not a global regularity
classification. The published sine convention is obtained from our
binary cosine convention by shifting the phase by $1/2$.

\subsection{Claims and boundaries}
\begin{center}
\begin{tabular}{@{}p{.24\textwidth}p{.69\textwidth}@{}}
\toprule
\textbf{Status} & \textbf{Content}\\
\midrule
Prior repository results & Leading square-root pressure, the atomic
Jordan block, and the leading bounded-detuning eigenvalue crossover.\\[2mm]
Proved here & Both next critical coefficients with an explicit
little-power-loss remainder; first normalized eigenvector corrections;
critical and detuned equilibrium selection; the right boundary layer.\\[2mm]
Derived consequences & An exact amplitude constant and the obstruction
to a pure half-integer-power expansion through order $|c|$.\\[2mm]
Not established & A complete transseries, a sharp next coefficient,
unbounded-detuning uniformity, a global phase theory, or formal verification.\\
\bottomrule
\end{tabular}
\end{center}

All model-specific analytic ingredients used below are proved or
reconstructed. We use the standard elementary framework of bounded
operators, Riesz projections, compact approximation, and measure entropy;
the particular finite-dimensional reduction and entropy argument are
spelled out. Numerical calculations are a separate diagnostic layer.

The pinned repository snapshot is
\begin{center}\small
\nolinkurl{1085b506d65e207a05b7e9c861bb1fe88432fe38}.
\end{center}
The inspected critical article has blob identifier
\nolinkurl{5539ac92693fc99fa35ded81bcf9a4486d37e501}.
The cited primary papers and targeted repository searches establish a
specific local research gap, not worldwide priority. Independent review
is still needed. No repository files were changed in preparing this work.

\section{Definitions and main results}
\label{sec:main}

\subsection{Normalization}
Fix $b\ge2$ and let $T_bx=bx\pmod1$ on $\TT=\RR/\ZZ$. Write
\begin{equation}
 m_b(x)=\left|\frac1b\sum_{r=0}^{b-1}e^{2\pi i r x}\right|
       =\left|\frac{\sin\pi bx}{b\sin\pi x}\right|,
 \label{eq:mask}
\end{equation}
with removable values at integers. For $s>0$ and phase $c$, define
\begin{equation}
 (\Lop_{s,c}f)(x)=\sum_{k=0}^{b-1}
 m_b\left(\frac{x+k}{b}-c\right)^s
 f\left(\frac{x+k}{b}\right).
 \label{eq:operator}
\end{equation}
There is no factor $1/b$ in this operator. Let $\mathcal D_{b,n}$ be the
$b^n$ equal closed cells of $[0,1]$. Put
\begin{align}
 F_{s,n,c}(x)&=\prod_{j=0}^{n-1}m_b(b^jx-c)^s,\\
 Z_{s,n}(c)&=\sum_{I\in\mathcal D_{b,n}}\sup_{x\in I}F_{s,n,c}(x),\\
 P_{b,s}(c)&=\lim_{n\to\infty}\frac1n\log Z_{s,n}(c).
 \label{eq:pressure}
\end{align}
The usual subdivision inequality makes $Z_{s,n}$ submultiplicative;
in the neighborhoods used here we will also identify the limit directly
with a simple eigenvalue. All logarithms are natural. Throughout,
\begin{equation}
 \delta=|c|,\quad t=\sqrt\delta,\quad d=b-1,\quad
 L=\log b,\quad a=2L,\quad \kappa_b=\sqrt{ad}.
 \label{eq:constants}
\end{equation}
Euler's constant is $\gamma$, and $\Leb$ denotes normalized Lebesgue
measure on $\TT$. The binary pressure order is
\begin{equation}
 P(q\log\cos^2\pi(x-c))=P_{2,2q}(c).
 \label{eq:binary}
\end{equation}
Thus our critical exponent $s=1$ corresponds to $q=1/2$.

\subsection{The complete expansion through the linear scale}
\begin{theorem}[Logarithmic critical correction]
\label{thm:pressure}
For every integer $b\ge2$ and every $0<\eta<1/2$, as $c\to0$,
\begin{equation}
 \boxed{\begin{aligned}
 P_{b,1}(c)={}&\kappa_b\sqrt\delta-d\delta\log(1/\delta)
                  +B_b\delta+O_{b,\eta}(\delta^{3/2-\eta}),\\
 B_b={}&b\log b+(b-1)(\gamma-1).
 \end{aligned}}
 \label{eq:main-pressure}
\end{equation}
For small nonzero $c$, the two eigenvalues near one satisfy
\begin{equation}
 \lambda_\pm(c)=1\pm\kappa_b\sqrt\delta-d\delta\log(1/\delta)
 +C_b\delta+O_{b,\eta}(\delta^{3/2-\eta}),
 \label{eq:eigenvalues-refined}
\end{equation}
where $C_b=(2b-1)\log b+(b-1)(\gamma-1)$.
The remaining spectrum on a suitable H\"older space stays in a fixed
disk of radius less than one, and $P_{b,1}(c)=\log\lambda_+(c)$.
\end{theorem}

For example,
\begin{equation}
\begin{split}
 P_{2,1}(c)={}&\sqrt{2\log2}\sqrt\delta-\delta\log(1/\delta)\\
 &+(2\log2+\gamma-1)\delta+O_\eta(\delta^{3/2-\eta}).
\end{split}
\label{eq:binary-expansion}
\end{equation}
Here $B_2\approx0.9635100260$. A polynomial in $\sqrt\delta$ cannot
represent the expansion through order $\delta$: the nonzero
$\delta\log\delta$ term must be retained. This is an obstruction to a
pure Puiseux expansion, not a proof of a complete logarithmic transseries.

\subsection{Eigenvectors, boundary layers, and selected measures}
Define, for $0<x<1$,
\begin{equation}
 H(x)=\frac{\sin\pi x}{\pi},\qquad
 G(x)=H(x)\left[-\psi(x)-\psi(1-x)+2\log\frac2\pi\right],
 \label{eq:HG}
\end{equation}
where $\psi=\Gamma'/\Gamma$. Extend periodically with $H(0)=0$ and
$G(0)=1$. For $f\in C^\alpha$, $0<\alpha<1$, set
\begin{equation}
 \ell(f)=f(0),\qquad
 \Xi(f)=\int_0^1\frac{f(x)-f(0)}{H(x)}\dd x.
 \label{eq:functionals}
\end{equation}
The integral is absolutely convergent. The functional $\Xi$ is bounded
on $C^\alpha$ but is not asserted to be a finite signed measure.

\begin{theorem}[Critical eigenvectors and half-and-half selection]
\label{thm:selection-critical}
Fix $0<\alpha<1/2$. For sufficiently small $c\ne0$ there are a strictly
positive right eigenfunction $h_c$ and a positive left eigenprobability
$\nu_c$ for $\lambda_+(c)$, uniquely normalized by
$\Xi(h_c)=1$ and $\nu_c(1)=1$. As $c\to0$,
\begin{align}
 h_c&=H+\frac{\kappa_b}{a}\sqrt\delta\,G+o(\sqrt\delta)
 &&\text{in }C^\alpha,\label{eq:right-expansion}\\
 \nu_c&=\ell+\frac{\kappa_b}{a}\sqrt\delta\,\Xi+o(\sqrt\delta)
 &&\text{in }(C^\alpha)^*.\label{eq:left-expansion}
\end{align}
The right boundary layer is
\begin{equation}
 \frac{h_c(\sqrt\delta\,z)}{\sqrt\delta}
 \longrightarrow |z|+\frac{\kappa_b}{a},
 \label{eq:boundary-critical}
\end{equation}
uniformly for $z$ in each compact subset of $\RR$, with the argument
interpreted on the circle. The probability
\begin{equation}
 \mu_c(f)=\frac{\nu_c(h_cf)}{\nu_c(h_c)}
 \label{eq:mu}
\end{equation}
is the unique equilibrium measure for $\log m_b(x-c)$, and
\begin{equation}
 \boxed{\displaystyle\mu_c\ \Longrightarrow\
             \tfrac12\delta_0+\tfrac12\Leb.}
 \label{eq:half-selection}
\end{equation}
The arrow means weak convergence of probabilities, not convergence in
total variation.
\end{theorem}

\begin{theorem}[Bounded-detuning selection]
\label{thm:selection-detuned}
Let $c=\sigma t^2$, $s=1+ut$, where $\sigma\in\{-1,1\}$ and $t\downarrow0$.
For $u$ in any fixed compact subset of $\RR$, the unique equilibrium
probabilities of $s\log m_b(x-c)$ satisfy, uniformly against each
continuous test function,
\begin{equation}
 \boxed{\displaystyle
 \mu_{1+ut,\sigma t^2}\Longrightarrow
 w_b(u)\delta_0+(1-w_b(u))\Leb,\qquad
 w_b(u)=\frac12\left(1+
 \frac{uL}{\sqrt{u^2L^2+4\kappa_b^2}}\right).}
 \label{eq:detuned-selection}
\end{equation}
At $(s,c)=(1,0)$ the entire set of equilibrium probabilities is exactly
$\{v\delta_0+(1-v)\Leb:0\le v\le1\}$. Thus the phase perturbation and its
detuning select particular elements of a pre-existing equilibrium segment.
\end{theorem}
The limiting weights tend algebraically to one or zero as
$u\to+\infty$ or $u\to-\infty$. The theorem does not make its error
uniform along an unbounded function $u=u(t)$.

\section{Reconstructing the critical spectral foundation}
\label{sec:jordan}
The Jordan mechanism in this section is a predecessor result
\cite{CriticalRepo}. We reconstruct it to fix the normalization of $G$
and the dual coordinates, which is essential for the next coefficients.
Our $G$ differs from the predecessor's unnormalized digamma companion by
a constant multiple of $H$.

\subsection{An exact Jordan chain and its dual}
Write $\Lop=\Lop_{1,0}$. The mask identity gives
\begin{equation}
 (\Lop f)(x)=\frac{H(x)}b\sum_{k=0}^{b-1}
             \frac{f((x+k)/b)}{H((x+k)/b)},\qquad 0<x<1.
 \label{eq:conjugacy}
\end{equation}
The removable endpoint values are understood by continuity.

\begin{lemma}[Normalized Jordan coordinates]
\label{lem:jordan-coordinates}
The functions and functionals in \cref{eq:HG,eq:functionals} satisfy
\begin{gather}
 \Lop H=H,\qquad \Lop G=G+aH,\qquad
 \ell\Lop=\ell,\qquad \Xi\Lop=\Xi+a\ell,\label{eq:jordan-relations}\\
 \ell(G)=1,\quad \ell(H)=0,\quad \Xi(H)=1,\quad \Xi(G)=0.
 \label{eq:normalization}
\end{gather}
Both $G$ and $H$ are periodic Lipschitz functions. Consequently
\begin{equation}
 \Proj f=G\ell(f)+H\Xi(f),\qquad
 \mathcal N f=aH\ell(f)
 \label{eq:PN}
\end{equation}
define a bounded rank-two projection and a nonzero nilpotent operator
with $\mathcal N^2=0$ and
$\Lop\Proj=\Proj\Lop=\Proj+\mathcal N$.
\end{lemma}
\begin{proof}
The expansions $\psi(x)=-x^{-1}-\gamma+O(x)$ and
$\psi(1-x)=-\gamma+O(x)$ give $G(0)=1$ with a finite one-sided derivative;
reflection handles the other endpoint. Formula \cref{eq:conjugacy}
immediately fixes $H$. The digamma multiplication identity
\[
 \sum_{k=0}^{b-1}\psi\left(\frac{x+k}{b}\right)
       =b\psi(x)-b\log b
\]
and its reflected version show that $\Lop G=G+2\log b\,H$.
The multiplication identity follows by logarithmic differentiation of
the Gamma multiplication formula, or directly by grouping the convergent
series for differences of digamma functions into residue classes.

For the only nontrivial normalization, integrate on
$[\varepsilon,1-\varepsilon]$. Since $\psi=(\log\Gamma)'$,
\begin{align*}
 \int_\varepsilon^{1-\varepsilon}
 [-\psi(x)-\psi(1-x)]\dd x
 &=2\log\Gamma(\varepsilon)-2\log\Gamma(1-\varepsilon),\\
 \int_\varepsilon^{1-\varepsilon}\frac{\dd x}{H(x)}
 &=2\log\cot(\pi\varepsilon/2).
\end{align*}
Using $\log\Gamma(\varepsilon)=-\log\varepsilon+O(\varepsilon)$,
the constant $2\log(2/\pi)$ in $G/H$ cancels the remaining finite part.
This proves $\Xi(G)=0$.

If $f(0)=0$, \cref{eq:conjugacy} and a change of variables give
$\Xi(\Lop f)=\Xi(f)$; absolute integrability follows from
$|f(x)|\le C\dist(x,\ZZ)^\alpha$. For arbitrary $f$, subtract $f(0)G$
and use the already proved formula for $\Lop G$. This proves the dual
identity. The projection and commutation assertions are now direct
consequences of the four pairings in \cref{eq:normalization}.
\end{proof}

\subsection{Weighted quadrature and the complementary spectrum}
\begin{lemma}[Uniform iterate estimate]
\label{lem:quadrature}
For $0<\alpha<1$, $N=b^n$, and $f\in C^\alpha$, one has
\begin{equation}
 \left\|\Lop^n f-a n H f(0)-G f(0)-H\Xi(f)\right\|_\infty
 \le C_{b,\alpha}(n+1)b^{-\alpha n}\|f\|_\alpha.
 \label{eq:C0-remainder}
\end{equation}
\end{lemma}
\begin{proof}
Iterating the telescoping quotient in \cref{eq:conjugacy} yields
\begin{equation}
 \Lop^n f(x)=\frac{H(x)}N\sum_{k=0}^{N-1}
             \frac{f((x+k)/N)}{H((x+k)/N)}.
 \label{eq:iterate}
\end{equation}
First let $f=1$. Subtract the two endpoint poles from the cosecant:
\[
 v(y)=\csc\pi y-\frac1{\pi y}-\frac1{\pi(1-y)}.
\]
This extends smoothly to $[0,1]$, and direct integration gives
$\int_0^1v=(2/\pi)\log(2/\pi)$. The two pole sums are evaluated using
\[
 \sum_{k=0}^{N-1}\frac1{x+k}=\psi(N+x)-\psi(x).
\]
Since $\psi(N+x)=\log N+O(N^{-1})$ uniformly for $0\le x\le1$, and
$H(x)$ removes the apparent endpoint divergences, the resulting expression
is $2\log N\,H(x)+G(x)+O(N^{-1})$, uniformly in $x$.

Now put $g=f-f(0)$. The function $g/H$ is integrable, although generally
unbounded. Compare the shifted sample in each cell of length $1/N$ with
its integral. For the first cell, the weighted sample is bounded by
\[
 \frac{H(x)}N\frac{|g(x/N)|}{H(x/N)}
 \le C\|f\|_\alpha N^{-\alpha}H(x)x^{\alpha-1}
 \le C'\|f\|_\alpha N^{-\alpha};
\]
the last cell is analogous. The endpoint integrals have the same bound.
For an interior cell at distance comparable to $k/N$ from an endpoint,
the variation of the numerator contributes at most
$C\|f\|_\alpha N^{-\alpha}/k$ after integration; the variation of $1/H$
contributes at most $C\|f\|_\alpha N^{-\alpha}k^{\alpha-2}$.
The same estimates hold with $N-k$ near the other endpoint.
Their sums are $O(N^{-\alpha}\log N)$ and $O(N^{-\alpha})$, respectively.
Thus the weighted sum converges uniformly to $H(x)\Xi(f)$ at the stated
rate. The case $n=0$ is absorbed by increasing the constant.
\end{proof}

\begin{proposition}[Spectral separation]
\label{prop:gap}
On $C^\alpha$, the spectral subspace of $\Lop$ at eigenvalue one is
exactly $\operatorname{span}\{G,H\}$, with its size-two Jordan block.
All remaining spectrum has modulus at most $b^{-\alpha}$. In particular,
for any $b^{-\alpha}<\theta<1$ there is $K_\theta<\infty$ such that
\begin{equation}
 \|\Lop^n-\Proj-n\mathcal N\|_{\alpha\to\alpha}
 \le K_\theta\theta^n\qquad(n\ge1).
 \label{eq:gap}
\end{equation}
\end{proposition}
\begin{proof}
For each fixed $n$, the inverse-branch formula and the product inequality
for H\"older seminorms give a Lasota--Yorke estimate
\begin{equation}
 [\Lop^n f]_\alpha\le
 C b^{-\alpha n}\|\Lop^n1\|_\infty[f]_\alpha+C_n\|f\|_\infty.
 \label{eq:LY}
\end{equation}
Here $C$ is independent of $n$. On arcs crossing the coordinate cut,
reindex the inverse branches and split at the cut; this changes only the
fixed H\"older constant. The coefficients are Lipschitz for each $n$.

To see the compactness consequence explicitly, approximate the identity
by periodic piecewise-linear interpolation $Q_M$ on a uniform mesh.
These finite-rank maps have uniformly bounded $C^\alpha$ norms and
$\|I-Q_M\|_{\alpha\to0}=O(M^{-\alpha})$. Apply \cref{eq:LY} to
$(I-Q_M)f$ and let $M\to\infty$. The essential norm of $\Lop^n$ is at
most $C' b^{-\alpha n}\|\Lop^n1\|_\infty$.
By \cref{lem:quadrature}, the last factor is $O(n+1)$, so
$\rhoess(\Lop)\le b^{-\alpha}$.

Outside that disk, spectrum consists of isolated eigenvalues of finite
algebraic multiplicity. If $\Lop f=\lambda f$ and $\lambda\ne1$, the dual
relations imply $\ell(f)=\Xi(f)=0$. Formula \cref{eq:C0-remainder}
then excludes $|\lambda|>b^{-\alpha}$. A generalized eigenvector at one,
after subtracting its $\Proj$ component, also lies in the common kernel
of the two functionals. Its iterates are a polynomial in $n$, whereas
\cref{eq:C0-remainder} makes them tend to zero. That polynomial must
vanish identically. The entire generalized eigenspace is therefore the
range of $\Proj$. On its complement the spectral-radius formula gives
\cref{eq:gap}, since the central block has exact powers
$\Proj+n\mathcal N$.
\end{proof}

\section{A perturbation reduction with a controlled chart error}
\label{sec:compression}

\subsection{The strength of the phase perturbation}
For the remainder of the critical calculation, let
$E_c=\Lop_{1,c}-\Lop_{1,0}$. The absolute trigonometric mask is Lipschitz,
so its translated difference has uniform norm $O(\delta)$ and bounded
Lipschitz seminorm. Splitting a H\"older quotient into distances below
and above $\delta$ proves
\begin{equation}
 \|E_c\|_{\alpha\to\alpha}=O(\delta^{1-\alpha})
 =O(\varepsilon_c),\qquad \varepsilon_c=\delta^{1-\alpha}.
 \label{eq:translation-bound}
\end{equation}
The finite inverse-branch sum is bounded on $C^\alpha$, so the weight
estimate indeed gives the operator estimate. It does not assert norm
differentiability of translation at a cusp.

\subsection{Graph coordinates around an isolated spectral cluster}
\begin{lemma}[Second-order compression error]
\label{lem:compression}
Let $A_0$ be a bounded operator with an isolated finite-dimensional
spectral subspace $V$ and commuting projection $P$. Put $Q=I-P$.
For sufficiently small bounded $E$, the corresponding spectral subspace
of $A_0+E$ is the graph of a map $K:V\to QX$ with
$\|K\|=O(\|E\|)$. In the graph coordinate $v\mapsto v+Kv$, the
compressed operator is
\begin{equation}
 A_0|_V+PE|_V+O(\|E\|^2).
 \label{eq:compression}
\end{equation}
The constants are uniform in a continuous family with a common separating
contour and uniformly bounded resolvents.
\end{lemma}
\begin{proof}
A resolvent Neumann series on the separating contour shows that the
Riesz projection $P_E$ satisfies $\|P_E-P\|=O(\|E\|)$.
The restriction of $P$ to $\operatorname{ran}P_E$ is invertible onto $V$,
so that range is a graph, with $K=O(\|E\|)$.
Projection back to $V$ gives the exact compressed expression
\[
 P(A_0+E)(I+K)|_V=A_0|_V+PE|_V+PEK.
\]
The term $PA_0K$ vanishes because $QX$ is $A_0$-invariant. The final term
has norm $O(\|E\|^2)$. This proves the claim without selecting individual
eigenvectors inside a possibly defective cluster.
\end{proof}

In the ordered basis $(G,H)$, \cref{lem:compression} and
\cref{prop:gap} give the two-dimensional matrix
\begin{equation}
 A_c=\begin{pmatrix}
 1+\ell(E_cG)&\ell(E_cH)\\
 a+\Xi(E_cG)&1+\Xi(E_cH)
 \end{pmatrix}+O(\delta^{2-2\alpha}).
 \label{eq:central-matrix}
\end{equation}
The error is an entrywise estimate in a fixed basis. It is not discarded
without examining its effect on the square-root splitting: an
$O(\delta^{2-2\alpha})$ error in the small upper-right entry can shift
an eigenvalue by $O(\delta^{3/2-2\alpha})$. Choosing $\alpha<1/4$ is what
makes this $o(\delta)$ and preserves the linear coefficient.

\section{Exact branch identities and the finite-part logarithm}
\label{sec:response}
Reflection conjugates $\Lop_{1,c}$ to $\Lop_{1,-c}$, so eigenvalues and
pressure are even. We therefore take $0<c<1/(4b)$ in this section and
restore $\delta=|c|$ at the end.

\subsection{A signed cotangent calculation}
\begin{lemma}[Exact response of the vanishing eigenfunction]
\label{lem:exactH}
For $d=b-1$,
\begin{equation}
 (E_cH)(0)=\frac{\sin\pi dc}{\pi}.
 \label{eq:atomH}
\end{equation}
On $bc\le x\le1$,
\begin{equation}
 (\Lop_{1,c}H)(x)=\frac{\sin\pi(x-dc)}\pi.
 \label{eq:bulkH}
\end{equation}
On $0\le x\le bc$,
\begin{equation}
 (\Lop_{1,c}H)(x)=
 2m_b(x/b-c)H(x/b)-\frac{\sin\pi(x-dc)}\pi.
 \label{eq:innerH}
\end{equation}
\end{lemma}
\begin{proof}
Set $y_k=(x+k)/b$. In the numerator absolute value, absorb the factor
$(-1)^k$ from $\sin\pi b(y_k-c)$, leaving $|\sin\pi(x-bc)|$.
Introduce the auxiliary signed sum obtained from this expression by
omitting its two absolute values. After multiplication by $H(y_k)$ it is
\[
 \frac{\sin\pi(x-bc)}{b\pi}
 \sum_{k=0}^{b-1}\frac{\sin\pi y_k}{\sin\pi(y_k-c)}.
\]
Use
$\sin\pi y_k/\sin\pi(y_k-c)=\cos\pi c+
\sin\pi c\cot\pi(y_k-c)$ and
\[
 \sum_{k=0}^{b-1}\cot\pi\left(\frac{x+k}{b}-c\right)
 =b\cot\pi(x-bc).
\]
The sum becomes $\sin\pi(x-dc)/\pi$. When $x\ge bc$, both the common
numerator and all denominators have the sign yielding the absolute-value
sum, so this proves \cref{eq:bulkH}. When $0<x<bc$, the branch $k=0$
has negative numerator and denominator, whereas the other denominators
remain positive. Restoring absolute values reverses the other branch
signs. The total is twice the zeroth-branch contribution minus the
signed sum, giving \cref{eq:innerH}. Continuity at $x=0$, where
$H(0)=0$, gives \cref{eq:atomH}. Apparent cotangent singularities are
removed by their sine prefactor.
\end{proof}

\subsection{The three matrix entries that determine the answer}
\begin{proposition}[Explicit response coefficients]
\label{prop:entries}
As $c\downarrow0$,
\begin{align}
 \ell(E_cH)&=dc+O(c^3),\label{eq:entry12}\\
 \ell(E_cG)&=r_bc+O(c^2),\label{eq:entry11}\\
 \Xi(E_cH)&=-2dc\log(1/c)+k_bc+O(c^2),\label{eq:entry22}
\end{align}
where
\begin{equation}
 r_b=2d\left(\gamma+\log\frac2\pi\right)+2b\log b,
 \qquad k_b=2d\left(\log\frac{\pi b}{2}-1\right).
 \label{eq:rk}
\end{equation}
For the remaining entry, the bound
$\Xi(E_cG)=O(c^{1-\alpha})$ suffices.
\end{proposition}
\begin{proof}
The first line is \cref{eq:atomH}. Evaluating $\Lop_{1,c}G$ at zero,
the branch at zero changes by $O(c^2)$ since $m_b(-c)=1+O(c^2)$.
For $1\le j<b$,
\[
 m_b(j/b-c)=\frac{\pi c}{\sin(\pi j/b)}+O(c^2)
            =\frac{c}{H(j/b)}+O(c^2).
\]
It follows that the coefficient in \cref{eq:entry11} is
$\sum_{j=1}^{b-1}G(j/b)/H(j/b)$. The digamma multiplication identity gives
\[
 \sum_{j=1}^{b-1}\psi(j/b)=-d\gamma-b\log b,
\]
which evaluates the coefficient as $r_b$.

For the logarithmic entry, write $v=\pi dc$. By definition,
\begin{equation}
 \Xi(E_cH)=\int_0^1
 \frac{\Lop_{1,c}H(x)-H(x)-\sin(v)/\pi}{H(x)}\dd x.
 \label{eq:finitepart-start}
\end{equation}
This is an ordinary convergent integral: the constant subtracted in the
numerator is the value of $E_cH$ at zero. Split it at $x=bc$.
On the outer interval, \cref{eq:bulkH} gives
\[
 \frac{\Lop_{1,c}H-H-\sin(v)/\pi}{H}
 =\cos v-1-\sin v\,[\cot\pi x+\csc\pi x].
\]
Since $\cot\pi x+\csc\pi x=\cot(\pi x/2)$, its integral is exactly
\begin{equation}
 I_{\rm out}=(1-bc)(\cos v-1)
       +\frac{2\sin v}{\pi}\log\sin\frac{\pi bc}{2}.
 \label{eq:outer-exact}
\end{equation}
Thus
\begin{equation}
 I_{\rm out}=-2dc\log(1/c)+2dc\log(\pi b/2)+O(c^2).
 \label{eq:outer-expand}
\end{equation}

On the inner interval, \cref{eq:innerH} and Taylor expansion give,
uniformly for $0\le x\le bc$ after removing the value at $x=0$,
\begin{equation}
 \frac{\Lop_{1,c}H-H-\sin(v)/\pi}{H(x)}
     =-\frac{2d}{b}+O(c^2).
 \label{eq:inner-expand}
\end{equation}
For a precise uniform justification, $m_b(x/b-c)=1+O(c^2)$ there,
$H(x/b)=x/b+O(x^3)$, and
\[
 \sin\pi(x-dc)+\sin\pi dc=\pi x+O(c^2x).
\]
The numerator in \cref{eq:inner-expand} is consequently
$-2dx/b+O(c^2x)$, and division by $H(x)=x(1+O(c^2))$ is uniform down
to the removable endpoint. Hence
$I_{\rm in}=-2dc+O(c^3)$. Combining this with
\cref{eq:outer-expand} proves \cref{eq:entry22}.
The last bound follows from boundedness of $\Xi$ and
\cref{eq:translation-bound}.
\end{proof}

\begin{corollary}[An exact binary check]
\label{cor:binary-exact}
For $b=2$ and $0<c<1/8$,
\begin{equation}
 \Xi(E_cH)=-2c+(1-2c)(\cos\pi c-1)
          +\frac{2\sin\pi c}{\pi}\log\sin\pi c.
 \label{eq:binary-exact}
\end{equation}
\end{corollary}
\begin{proof}
For the binary mask, the inner formula simplifies by elementary
sine addition to $\Lop_{1,c}H(x)=\sin\pi c/\pi$ when $0\le x\le2c$.
The integrand in \cref{eq:finitepart-start} is therefore exactly $-1$
on that interval. Add its integral $-2c$ to \cref{eq:outer-exact}.
\end{proof}

The appearance of $\log c$ is now transparent. The shifted zeros feed
a nonzero value into a function which previously vanished at zero.
The dual coordinate removes that value but integrates the residual
against $1/H(x)\sim1/\dist(x,\ZZ)$. The outer region beginning at $bc$
therefore leaves a logarithmic finite part. No divergent expression has
been treated as a convergent integral.

\section{From the response matrix to the pressure expansion}
\label{sec:proof-pressure}

\subsection{The refined eigenvalue calculation}
\begin{proof}[Proof of the eigenvalue assertion in \cref{thm:pressure}]
Fix $0<\eta<1/2$ and choose $\alpha=\eta/2<1/4$. Write the matrix in
\cref{eq:central-matrix} as
\[
 A_c=\begin{pmatrix}1+e&u\cr a+g&1+f\end{pmatrix}.
\]
Its entries, including the graph error, satisfy
\begin{align}
 e&=r_bc+O(c^{2-2\alpha}),\qquad
 f=-2dc\log(1/c)+k_bc+O(c^{2-2\alpha}),\\
 u&=dc+O(c^{2-2\alpha}),\qquad
 g=O(c^{1-\alpha}).
 \label{eq:entry-errors}
\end{align}
The two eigenvalues are exactly
\begin{equation}
 1+\frac{e+f}{2}\ \pm\
 \sqrt{u(a+g)+\frac{(e-f)^2}{4}}.
 \label{eq:roots}
\end{equation}
The expression under the square root is
\begin{equation}
 ad c+O(c^{2-2\alpha})+O(c^{2-\alpha})
                 +O(c^2\log^2 c).
 \label{eq:discriminant}
\end{equation}
It is positive for sufficiently small $c>0$. Dividing the error by
$\sqrt c$ in the difference-of-square-roots identity gives
\begin{equation}
 \sqrt{u(a+g)+(e-f)^2/4}
   =\kappa_b\sqrt c+O(c^{3/2-\eta}).
 \label{eq:sqrt-error}
\end{equation}
Here $c^{3/2}\log^2c=O_\eta(c^{3/2-\eta})$.

The half-trace is
\begin{equation}
 \frac{e+f}{2}=-dc\log(1/c)+\frac{r_b+k_b}{2}c
                      +O(c^{2-2\alpha}).
\end{equation}
The constants simplify to
\begin{equation}
 \frac{r_b+k_b}{2}=d(\gamma-1)+(2b-1)\log b=C_b.
 \label{eq:Cconstant}
\end{equation}
Combining these equations proves \cref{eq:eigenvalues-refined} for
positive phases. Reflection supplies the negative phases. The rank-two
cluster is the only spectrum near one; the complementary spectrum stays
in a fixed smaller disk by resolvent continuity. The eigenvalues are
real, simple, and separated by $2\kappa_b\sqrt\delta+o(\sqrt\delta)$.
\end{proof}

\subsection{Positivity and identification with cylinder pressure}
\begin{lemma}[Positive spectral data and pressure]
\label{lem:positive-pressure}
For sufficiently small $c\ne0$, $\lambda_+(c)$ is the simple dominant
eigenvalue of $\Lop_{1,c}$, with a strictly positive right eigenfunction
and a positive left eigenmeasure. Its logarithm equals the pressure in
\cref{eq:pressure}.
\end{lemma}
\begin{proof}
Use the graph coordinate of \cref{lem:compression}. A right eigenvector
in the central plane can be written $(x_c,1)$, where the second row of
the eigenvalue equation gives
\begin{equation}
 x_c=\frac{\lambda_+(c)-1-f}{a+g}
     =\frac{\kappa_b}{a}\sqrt\delta+o(\sqrt\delta)>0.
 \label{eq:x-leading}
\end{equation}
Lifting to the full space gives
\begin{equation}
 h_c=H+x_cG+O(\delta^{1-\alpha})
 \label{eq:h-lift}
\end{equation}
in $C^\alpha$. Choose a small fixed neighborhood of zero on which
$G\ge1/2$. There $H\ge0$ and $x_c$ is a positive constant times
$\sqrt\delta$, while the error is $o(\sqrt\delta)$ when $\alpha<1/2$.
Outside that neighborhood, $H$ is bounded below by a positive constant.
This proves $h_c>0$ everywhere for sufficiently small $c\ne0$.

The positive operator $\Lop_{1,c}$ has the norm limit
$\lambda_+^{-n}\Lop_{1,c}^n\to\Proj_+$ on $C^\alpha$ for each fixed
$c\ne0$, since all its other spectrum is strictly smaller in modulus.
Thus $\Proj_+$ is positive. Write $\Proj_+f=h_c\ell_c(f)$, with
$\ell_c(h_c)=1$. Its positivity makes $\ell_c$ a positive functional.
It is bounded for the uniform norm, because
$|\ell_c(f)|\le\ell_c(1)\|f\|_\infty$ for real $f$.
Its extension to continuous functions is a positive finite measure.
Normalize this measure to obtain $\nu_c(1)=1$.

A change of variables gives
\begin{equation}
 \int_0^1\Lop_{1,c}^n1\dd x=b^n\int_0^1 F_{1,n,c}\dd x
          \le Z_{1,n}(c).
 \label{eq:pressure-lower}
\end{equation}
The left side has positive leading coefficient times $\lambda_+^n$, so
it supplies the lower pressure bound.
For the upper bound, fix $c\ne0$ and replace the weight $m_b(x-c)$ by
$m_b(x-c)+\epsilon$. Its operator converges in norm to $\Lop_{1,c}$
as $\epsilon\downarrow0$. The simple leading eigenvalue and positive
eigenfunction persist, say as $\lambda_\epsilon$ and $h_\epsilon$.
The logarithm of this regularized weight is H\"older continuous and
bounded. On each $b$-adic cylinder, inverse-branch contraction bounds
the variation of its $n$-step logarithmic weight by a geometric series,
independently of $n$. Cylinder suprema are consequently comparable,
with an $n$-independent factor, to the weighted inverse-branch sum at any
fixed target point. That sum is $\Lop_\epsilon^n1$ and has exponential
rate $\lambda_\epsilon$. Hence the regularized cylinder pressure is
$\log\lambda_\epsilon$. Since the original products are pointwise no
larger, their upper exponential rate is at most
$\log\lambda_\epsilon$. Letting $\epsilon\downarrow0$ proves the claim.
This argument needs no bounded logarithm for the original weight.
\end{proof}

\begin{proof}[Completion of \cref{thm:pressure}]
Expand the logarithm of \cref{eq:eigenvalues-refined} for the positive
branch. The quadratic term subtracts $\kappa_b^2\delta/2=d\log b\,\delta$.
All cross terms involving $\sqrt\delta$ and $\delta\log\delta$, and all
higher Taylor terms, are $O_\eta(\delta^{3/2-\eta})$.
Thus $B_b=C_b-d\log b$, which is exactly \cref{eq:main-pressure}.
At $c=0$, the telescoping product is $H(b^nx)/(b^nH(x))$. On a cell
at distance comparable to $k/b^n$ from an endpoint its supremum is
$O(1/k)$, while the two endpoint cells contribute $O(1)$. Thus
$1\le Z_{1,n}(0)\le C(n+1)$, proving $P_{b,1}(0)=0$. The expansion
therefore extends continuously to zero.
\end{proof}

\begin{remark}[What the remainder does and does not say]
For each fixed $\eta>0$, the remainder is smaller than $\delta$.
It is not asserted to be uniform as $\eta\downarrow0$, and no explicit
numerical bound for its constant is claimed. The calculation leaves
terms at the scale $\delta^{3/2}\log^2\delta$ unresolved. A generic
operator norm estimate is sufficient for both displayed coefficients,
but not for a sharp subsequent term.
\end{remark}

\section{Individual eigenvectors and the critical boundary layer}
\label{sec:eigenvectors}

\subsection{Two normalizations that remain meaningful at the limit}
Choose now any fixed $0<\alpha<1/2$. The leading eigenvalue expansion,
which does not require the refined remainder on this particular space,
is $\lambda_+=1+\kappa_b t+o(t)$. The graph error is
$O(\delta^{1-\alpha})=o(t)$. Since graph vectors lie in the kernel of
both $\ell$ and $\Xi$, the lift of $(x_c,1)$ has $\Xi(h_c)=1$ exactly.
Equation \cref{eq:x-leading} proves
\cref{eq:right-expansion}.

A left eigenvector of the central matrix can be written $(1,q_c)$.
Its first component equation gives
\begin{equation}
 q_c=\frac{\lambda_+(c)-1-e}{a+g}
          =\frac{\kappa_b}{a}t+o(t).
 \label{eq:q-leading}
\end{equation}
To lift this row to the full dual space, compose it with the coordinate
map from the perturbed cluster projection. That map equals
$(\ell,\Xi)+O(\delta^{1-\alpha})$ in operator norm. The resulting
functional is
$\ell+q_c\Xi+o(t)$. Because $\Xi(1)=0$, its value at one is $1+o(t)$;
normalizing that value to one proves \cref{eq:left-expansion}.
Positivity and uniqueness were established in \cref{lem:positive-pressure}.

It follows, in particular, that
\begin{equation}
 \nu_c\Longrightarrow\delta_0,\qquad
 h_c\longrightarrow H,\qquad
 \nu_c(h_c)=\frac{2\kappa_b}{a}t+o(t).
 \label{eq:normalization-collapse}
\end{equation}
The weak convergence follows first on H\"older test functions, then on
all continuous functions by uniform approximation and the fact that the
$\nu_c$ are probabilities.

\subsection{The boundary layer}
On a fixed compact set of real $z$, periodicity and the expansion of
sine give $H(tz)/t\to|z|$, whereas $G(tz)\to1$ uniformly.
The $C^\alpha$ remainder in \cref{eq:right-expansion} is $o(t)$ even in
uniform norm. This proves \cref{eq:boundary-critical}.

The scaling is not arbitrary: the bulk eigenfunction vanishes linearly
at the atom, while the positive lift at the atom is of order $t$.
Balancing these two terms forces the spatial scale $x\asymp t$.
This statement concerns the right eigenfunction; it does not assert
that the left eigenmeasure has a density with a comparable pointwise
profile.

\subsection{Why the weak left limit is not the equilibrium limit}
For any fixed $f\in C^\alpha$, multiply the two expansions and use
\begin{equation}
 \Xi(Hf)=\int_0^1f\dd x,
 \label{eq:key-pairing}
\end{equation}
since $H(0)=0$. One obtains
\begin{equation}
 \nu_c(h_cf)=\frac{\kappa_b}{a}t
        \left(f(0)+\int_0^1f\dd x\right)+o(t).
 \label{eq:pairing-critical}
\end{equation}
Divide by \cref{eq:normalization-collapse} to obtain the half-and-half
limit in \cref{eq:half-selection}. Uniform approximation extends it to
continuous test functions.

A naive multiplication of the limiting objects would give
$H\delta_0=0$, not a probability. Both the numerator and denominator in
\cref{eq:mu} vanish at order $t$. Their first corrections contain the
information that selects a nontrivial mixture. Also, the expansion of
$\nu_c$ is in $(C^\alpha)^*$, not in the total-variation norm; the
finite-part functional $\Xi$ must not be mistaken for a signed
probability correction.

\begin{corollary}[Exact leading amplitude divergence]
\label{cor:amplitude}
For each fixed sufficiently small $c\ne0$ there are $\theta_c<1$ and
$\mathcal A_b(c)>0$ such that
\begin{equation}
 b^n\int_0^1F_{1,n,c}\dd x
 =\mathcal A_b(c)e^{nP_{b,1}(c)}(1+O_c(\theta_c^n)),
 \label{eq:amplitude-n}
\end{equation}
and
\begin{equation}
 \mathcal A_b(c)\sim\frac{a}{\pi^2\kappa_b}|c|^{-1/2}.
 \label{eq:amplitude-c}
\end{equation}
\end{corollary}
\begin{proof}
The rank-one spectral projection is
$f\mapsto h_c\nu_c(f)/\nu_c(h_c)$.
Integrate its value at $f=1$ and use
$\int H=2/\pi^2$ and \cref{eq:normalization-collapse}.
The complementary spectrum gives the fixed-$c$ exponential error.
\end{proof}
The divergence order is already visible in the predecessor's spectral
projectors. Formula \cref{eq:amplitude-c} records its exact constant for
the normalization \cref{eq:operator}. The two successive limits in this
corollary are not a new uniform finite-size theorem.
% ed. (2026-09-29): the amplitude asymptotic is already in the predecessor.
\begin{ednote}
The asymptotic \cref{eq:amplitude-c} is not new: the predecessor
\cite{CriticalRepo} already states, for its two simple eigenprojections
$P_\pm(c)$, that
$\int_0^1P_\pm(c)1\dd x\sim\pm a_b/(\pi^2\kappa_b\sqrt{|c|})$ with
$a_b=2\log b$ (its display labelled \texttt{eq:critical-amplitudes}), and
$\mathcal A_b(c)$ here is exactly $\int_0^1P_+(c)1\dd x$. What this article
adds is the derivation from the normalized eigenvectors of
\cref{thm:selection-critical}. The same applies to the corresponding
phrases in the abstract and in the claims table of \cref{sec:intro}.
\end{ednote}

\section{Detuned eigenvectors and the selected mixture}
\label{sec:detuning}
The leading eigenvalue crossover used in this section was already proved
in \cite{CriticalRepo}. We give a short reconstruction that also produces
the eigenvectors needed for the new selection statement.

\subsection{An analytic atomic cluster}
For real $s$ close to one, put
$H_s=H^s$ and $\beta_s=b^{1-s}$. The exact coboundary gives
\begin{equation}
 \Lop_{s,0}H_s=\beta_s H_s,\qquad
 \ell\Lop_{s,0}=\ell.
 \label{eq:atomic-detuned}
\end{equation}
For fixed $\alpha<1$, the map $s\mapsto\Lop_{s,0}$ is analytic in a
complex neighborhood of one on $C^\alpha$. At a simple zero, the local
factor is $|x|^s$ times a nonvanishing analytic factor; the derivatives
$|x|^s\log^k|x|$ have the necessary H\"older bounds uniformly whenever
$\Re s>\alpha$. Away from the zeros analyticity is immediate. The same
reasoning applies to $H_s$.

Let $\Proj_s$ be the analytic rank-two cluster projection, define
$G_s=\Proj_s1$, and write its second coordinate as $\Xi_s$, so that
\begin{equation}
 \Proj_sf=G_s\ell(f)+H_s\Xi_s(f).
 \label{eq:Ps}
\end{equation}
Indeed $\ell\Proj_s=\ell$, because $\ell$ is a left eigenvector at one
inside the contour. Hence $G_s(0)=1$ and the kernel of $\ell$ in the
central plane is generated by $H_s$. The coordinates satisfy
\begin{equation}
 G_1=G,\quad \Xi_1=\Xi,\quad \Xi_s(H_s)=1,\quad \Xi_s(G_s)=0.
 \label{eq:coords-detuned}
\end{equation}
They depend analytically on $s$. In this basis the atomic central matrix
has the form
\begin{equation}
 A_{s,0}=\begin{pmatrix}1&0\\a_s&\beta_s\end{pmatrix},
 \qquad a_s=a+O(s-1).
 \label{eq:detuned-matrix0}
\end{equation}
No explicit formula for $a_s$ is needed for leading selection.

\subsection{Balanced coordinates and the two eigenvectors}
Set $c=\sigma t^2$ and $s=1+ut$. Fix $0<\alpha<1/2$ and let $u$ range
in a compact set. The phase perturbation has norm
\begin{equation}
 \|\Lop_{s,c}-\Lop_{s,0}\|_{\alpha\to\alpha}
       =O(t^{2-2\alpha})=o(t),
 \label{eq:detuned-small}
\end{equation}
uniformly. For $s<1$ use the local exponent $s-\alpha$; the extra factor
$t^{2ut}=\exp(2ut\log t)$ stays bounded and tends to one. For $s\ge1$
the uniform Lipschitz estimate suffices.

Evaluation at the atom supplies a sharper estimate for the small
upper-right entry:
\begin{equation}
 \ell\bigl((\Lop_{s,c}-\Lop_{s,0})H_s\bigr)
       =d|c|^s(1+O(t^2))=dt^2(1+o(1)).
 \label{eq:detuned-small-entry}
\end{equation}
To check this, each branch $j/b$, $1\le j<b$, contributes
$|c|^s(1+O(c))$, while the zeroth branch vanishes. The formula is uniform
for $s$ near one and for either sign of $c$.

Compress by the graph lemma and change coordinates by
$S_t=\operatorname{diag}(t,1)$. The square of the error in
\cref{eq:detuned-small} is $o(t^2)$, so it remains negligible even in
the small entry. Thus
\begin{equation}
 S_t^{-1}A_{s,c}S_t
 =I+t\begin{pmatrix}0&d\\a&-uL\end{pmatrix}+o(t).
 \label{eq:balanced}
\end{equation}
Define
\begin{equation}
 D(u)=\sqrt{u^2L^2/4+\kappa_b^2},\quad
 A(u)=D(u)-uL/2,\quad B(u)=D(u)+uL/2.
 \label{eq:AB}
\end{equation}
These are positive, with $A(u)B(u)=ad$ and $A(u)+B(u)=2D(u)$.
The dominant generator eigenvalue is $A(u)$, while its right balanced
vector is $(B(u)/a,1)$ and its left balanced row is proportional to
$(1,A(u)/a)$. Undoing the scaling and lifting back to the function space
gives, uniformly for bounded $u$,
\begin{align}
 \lambda_+(s,c)&=1+tA(u)+o(t),\label{eq:old-crossover}\\
 h_{s,c}&=H_s+\frac{B(u)}a tG_s+o(t),\label{eq:right-detuned}\\
 \nu_{s,c}&=\ell+\frac{A(u)}a t\Xi_s+o(t).
 \label{eq:left-detuned}
\end{align}
The normalizations are $\Xi_s(h_{s,c})=1$ and $\nu_{s,c}(1)=1$.
For the latter, $\Xi_s(1)=O(s-1)=O(t)$, so normalization changes no
first-order term. The argument for strict positivity in
\cref{lem:positive-pressure} applies uniformly on compact $u$ sets:
$B(u)$ has a positive minimum there and the graph error is $o(t)$.
The same regularization argument identifies the pressure with
$\log\lambda_+$.

\subsection{The mixture weights}
For each $f\in C^\alpha$,
$\Xi_s(H_sf)\to\Xi(Hf)=\int f\dd x$. Pairing
\cref{eq:right-detuned,eq:left-detuned} therefore gives
\begin{align}
 \nu_{s,c}(h_{s,c}f)
    &=\frac{t}{a}\left(B(u)f(0)+A(u)\int f\dd x\right)+o(t),\\
 \nu_{s,c}(h_{s,c})&=\frac{2D(u)}a t+o(t).
 \label{eq:detuned-pairings}
\end{align}
The denominator is uniformly bounded below by a positive multiple of
$t$ for bounded $u$. Dividing proves \cref{eq:detuned-selection}, since
$w_b(u)=B(u)/(2D(u))$. Uniform approximation extends the conclusion to
each continuous test function. This proves the weak-limit part of
\cref{thm:selection-detuned}, once the equilibrium interpretation is
established below.

The same expansion gives the detuned right boundary layer
\begin{equation}
 \frac{h_{1+ut,\sigma t^2}(tz)}t
       \longrightarrow |z|+\frac{B(u)}a,
 \label{eq:boundary-detuned}
\end{equation}
uniformly for $(u,z)$ in compact sets. In fact
$H(tz)^{1+ut}/t\to|z|$, including $z=0$, and $G_s(tz)\to1$.
Positive detuning increases the atomic lift and the selected atomic
weight; negative detuning favors the Lebesgue component. These statements
are about the bounded scaling window, not fixed-$s$ extrapolations.

\section{Equilibrium measures and the critical equilibrium segment}
\label{sec:equilibrium}

\subsection{A variational argument that allows zeros}
For a weight $w=m_b(\,\cdot-c)^s$, an equilibrium probability means a
$T_b$-invariant probability maximizing
\begin{equation}
 h_\mu(T_b)+\int\log w\dd\mu.
 \label{eq:free-energy}
\end{equation}
The value $-\infty$ is permitted for the integral. The potential is not
bounded below, so this interpretation should be justified rather than
attached automatically to a transfer eigenvector.

\begin{lemma}[Positive spectral data give the unique equilibrium]
\label{lem:equilibrium}
Suppose the transfer operator for a nonnegative continuous H\"older
weight has a strictly positive eigenfunction $h$, a positive finite left
eigenmeasure $\nu$, and a simple dominant eigenvalue $\lambda$ on
$C^\alpha$. Then $\mu=h\nu/\nu(h)$ is invariant and is the unique
equilibrium probability. Its free energy equals $\log\lambda$.
\end{lemma}
\begin{proof}
Define the normalized inverse-branch weights
\begin{equation}
 g(y)=\frac{w(y)h(y)}{\lambda h(T_by)}.
 \label{eq:g}
\end{equation}
They are nonnegative and sum to one over each fiber of $T_b$.
For continuous $f$,
$\Lop(h(f\circ T_b))=\lambda hf$, which, paired with $\nu$, proves
invariance of $\mu$. More generally, for continuous $F$ and $K$,
\[
 \int F(y)K(T_by)\dd\mu(y)
 =\int K(x)\sum_{T_by=x}g(y)F(y)\dd\mu(x).
\]
Thus the conditional distribution of a predecessor, given its image,
is the finite probability vector $g$.

For completeness, the entropy formula needed here is
\begin{equation}
 h_\rho(T_b)=\int\left[-\sum_{T_by=x}p_x(y)\log p_x(y)\right]\dd\rho(x),
 \label{eq:entropy}
\end{equation}
where $p_x$ is the conditional predecessor distribution of any invariant
probability $\rho$. The partition into the $b$ half-open digit intervals
is generating. In its stationary digit process, the conditional entropy
of the present digit given the entire future equals the entropy rate:
for finite futures it is the difference of two consecutive block
entropies, by stationarity, and these differences decrease to that rate.
The image point encodes the entire future, so conditioning on it gives
\cref{eq:entropy}. The half-open convention handles the countable
ambiguity of terminating expansions and does not change the argument.

For $\rho=\mu$, \cref{eq:entropy} is at most $\log b$. Therefore
$-\log g$ is $\mu$-integrable, with $0\log0=0$. Since $\log h$ is bounded,
\cref{eq:g} makes $\log w$ integrable as well. Its coboundary integrates
to zero, giving free energy $\log\lambda$.

For arbitrary invariant $\rho$ with finite potential integral,
\cref{eq:g,eq:entropy} express the free energy minus $\log\lambda$ as
\begin{equation}
 -\int\sum_{T_by=x}p_x(y)\log\frac{p_x(y)}{g(y)}\dd\rho(x)\le0.
 \label{eq:relative-entropy}
\end{equation}
The inequality is the finite-vector relative-entropy inequality, with
the value $+\infty$ when positive mass is assigned to $g=0$.
A measure with potential integral $-\infty$ cannot maximize.
Equality forces $p_x=g$ almost everywhere. For such a measure, the finite
measure $h^{-1}\rho$ is a left eigenmeasure of $\Lop$ at $\lambda$.
The adjoint eigenspace there is one-dimensional because the eigenvalue
is isolated and algebraically simple on $C^\alpha$. Thus the measure is
proportional to $\nu$, and normalization forces $\rho=\mu$.
\end{proof}

The hypotheses hold at all sufficiently small nonzero phases considered
in \cref{thm:selection-critical,thm:selection-detuned}. This completes
their equilibrium assertions without invoking a positivity theorem that
would fail at the defective atomic operator.

\subsection{What is being selected at the critical point?}
\begin{proposition}[The critical equilibrium segment]
\label{prop:segment}
For $s=1$ and $c=0$, the equilibrium probabilities are exactly
$v\delta_0+(1-v)\Leb$, $0\le v\le1$.
\end{proposition}
\begin{proof}
Let $\rho$ be ergodic and different from $\delta_0$. Invariance implies
that the countable union of preimages of zero has $\rho$-measure zero.
For almost every $x$, the telescoping relation is therefore valid:
\begin{equation}
 \sum_{j=0}^{n-1}\log m_b(T_b^jx)
 =\log H(T_b^nx)-\log H(x)-n\log b.
 \label{eq:ergodic-coboundary}
\end{equation}
There is some $\epsilon>0$ for which $\rho(H\ge\epsilon)>0$.
Ergodicity gives infinitely many returns of almost every point to this
set. Along those return times, the average in
\cref{eq:ergodic-coboundary} tends to $-\log b$.
Since $-\log m_b\ge0$, the ergodic theorem with truncation implies that
its averages tend to $+\infty$ if its integral is infinite, and to its
integral otherwise. The return subsequence excludes infinity. Hence
\begin{equation}
 \int\log m_b\dd\rho=-\log b.
 \label{eq:ergodic-potential}
\end{equation}
Every such ergodic measure has free energy
$h_\rho(T_b)-\log b\le0$. Equality requires entropy $\log b$.
In the stationary digit process, maximal entropy means that the
conditional law of each digit given its past is the uniform distribution;
the digits are therefore independent and uniform. The corresponding
circle measure is exactly $\Leb$.

The fixed-point measure $\delta_0$ has both entropy and potential
integral zero, so it also has free energy zero. Entropy is affine under
ergodic decomposition, and the potential integral decomposes as well.
Consequently a general invariant probability attains zero precisely
when almost all of its ergodic components are $\delta_0$ or $\Leb$.
Their mixtures have free energy zero, which is the maximum and agrees
with $P_{b,1}(0)=0$.
\end{proof}
This also proves the last assertion of \cref{thm:selection-detuned}.
The segment classification is supplied to identify the variational
meaning of the limit; the phase-selected weights and boundary behavior
are the new conclusions targeted here.

\section{Numerical diagnostics and reproducibility}
\label{sec:numerics}

\subsection{Independent checks of the response identities}
The accompanying \code{code/verify.py} records its actual output in
\code{data/verification.json}. At 70-decimal-digit working precision,
20 choices of base and point test both Jordan identities, and eight
base--phase pairs test both the exact atom response and the bulk sine
identity. The largest absolute residuals observed in these two groups
are approximately $1.4\times10^{-70}$ and $1.7\times10^{-71}$.
These are high-precision numerical checks of independently proved
identities, not interval-arithmetic certificates.

The script also evaluates the finite-part integral directly after
splitting at $bc$. At $c=10^{-7}$, it gives the following comparisons:
\begin{center}\small
\begin{tabular}{@{}rrrrr@{}}
\toprule
$b$ & $\ell(E_cG)/c$ & $r_b$ & $\Xi(E_cH)/c+2d\log(1/c)$ & $k_b$\\
\midrule
2 & 3.023854148 & 3.023854641 & 0.289459278 & 0.289459772\\
3 & 7.094204255 & 7.094205570 & 2.200778002 & 2.200779976\\
5 & 17.099438853 & 17.099442801 & 8.488157046 & 8.488164942\\
\bottomrule
\end{tabular}
\end{center}
The residual differences in this table are mainly finite-$c$ corrections,
not claims of numerical error in the limiting constants. For the binary
closed formula \cref{eq:binary-exact}, two direct integrals with an endpoint
cutoff $10^{-30}$ differ from the exact expression by about $10^{-30}$,
as expected from omitting a bounded endpoint integrand.

\subsection{Pressure: a deliberately demanding observable}
For the continuum pressure define the logarithm-corrected quotient
\begin{equation}
 R_b(c)=\frac{P_{b,1}(c)-\kappa_b\sqrt{|c|}
                      +(b-1)|c|\log(1/|c|)}{|c|}.
 \label{eq:Rnumeric}
\end{equation}
The theorem predicts $R_b(c)\to B_b$.
The code discretizes the transfer operator by nonnegative periodic
piecewise-linear interpolation at the inverse branches and computes its
two eigenvalues with largest real parts. The following are floating-point
collocation values, not certified continuum values:
\begin{center}\small\begin{tabular}{@{}rrrrr@{}}
\toprule
$b$ & $c$ & $P_{b,1}(c)/\sqrt{|c|}$ & $R_b(c)$ & $B_b$\\
\midrule
2 & $10^{-3}$ & 1.02347259 & 2.039826 & 0.963510\\
2 & $10^{-4}$ & 1.10259410 & 1.728748 & 0.963510\\
2 & $10^{-5}$ & 1.14540630 & 1.392461 & 0.963510\\
2 & $10^{-6}$ & 1.16477281 & 1.178294 & 0.963510\\
3 & $10^{-3}$ & 1.80542336 & 4.617369 & 2.450268\\
3 & $10^{-4}$ & 1.95254224 & 4.045490 & 2.450268\\
3 & $10^{-5}$ & 2.03409965 & 3.358224 & 2.450268\\
3 & $10^{-6}$ & 2.07156876 & 2.905634 & 2.450268\\
\bottomrule
\end{tabular}
\end{center}
The leading constants are $\kappa_2\approx1.17741002$ and
$\kappa_3\approx2.09629415$. The displayed approach of $R_b(c)$ to $B_b$
is slow. That is compatible with the remainder and with additional
$|c|^{3/2}\log^2|c|$-scale contributions; the table does not establish
a coefficient for any such term.

The nominal meshes are $65\,536$ points at $c=10^{-3}$,
$262\,144$ at $10^{-4}$ and $10^{-5}$, and $1\,048\,576$ at $10^{-6}$.
They are rounded upward to a multiple of the base. A doubled-grid check
for $b=2$, $c=10^{-6}$ changes the pressure from
$0.00116477280617$ to $0.00116476829537$.
The change is approximately $4.51\times10^{-9}$ in pressure but about
$0.00451$ after division by $c$ in \cref{eq:Rnumeric}. This illustrates
why a small eigenvector residual alone does not certify the continuum
coefficient. Both meshes still give values above the predicted limiting
constant, consistent with nonnegligible finite-phase corrections.

\subsection{Selection: an observable distinguishing the two components}
The test function $\cos2\pi x$ has value one at the atom and integral
zero against $\Leb$. Hence its predicted limiting expectation is exactly
$w_b(u)$. Computing left and right eigenvectors of the same positive
collocation matrix and normalizing their entrywise product gives:
\begin{center}\small\begin{tabular}{@{}rrrr@{}}
\toprule
$u$ & $c$ & numerical $\int\cos(2\pi x)\,d\mu$ & predicted limit\\
\midrule
-2 & $10^{-4}$ & 0.26869569 & 0.24633959\\
-2 & $10^{-5}$ & 0.25711076 & 0.24633959\\
0 & $10^{-4}$ & 0.51963407 & 0.50000000\\
0 & $10^{-5}$ & 0.50961315 & 0.50000000\\
2 & $10^{-4}$ & 0.77612951 & 0.75366041\\
2 & $10^{-5}$ & 0.76440084 & 0.75366041\\
\bottomrule
\end{tabular}
\end{center}
All these rows use $b=2$ and $262\,144$ mesh points. The recorded output
also includes the sine moment, matrix residuals, and the discrepancy
between the eigenvalues obtained from left and right calculations.
The agreement is qualitative and improving across the two phases;
it is not a rigorous weak-convergence certificate. No numerical result
is used as a premise in any theorem above.

\subsection{How to reproduce the package}
With the versions recorded in \code{requirements.txt}, run
\begin{verbatim}
python -m pip install -r requirements.txt
OPENBLAS_NUM_THREADS=1 python code/verify.py
\end{verbatim}
The environment-variable syntax shown is for a POSIX shell; it only
limits numerical threading. On other shells set it using their syntax
or omit it. The script records library versions and computation residuals.
% ed. (2026-09-29): the default output directory is now recomputed/; double
% hyphens are set with -{}- (\code set them as en dashes in this font).
To skip collocation and retain only the special-function diagnostics,
use \texttt{python code/verify.py -{}-identities-only}. By default the
outputs go to \code{recomputed/}; the option \texttt{-{}-output-dir}
selects another directory, and only \texttt{-{}-output-dir data}
overwrites the shipped receipt and tables. The PDF is rebuilt by the supplied
\code{Makefile}, or by three successive runs of pdfLaTeX. No network
access is needed once the dependencies are installed.

\section{Further research questions and concrete next targets}
\label{sec:questions}
The questions below are proposed continuations of this article. They are
not all claimed to be previously published open problems. The first,
fourth, and fifth also refine explicitly identified boundaries in the
critical predecessor.

\begin{question}[The next critical logarithmic polynomial]
Does the next pressure term have the form
$|c|^{3/2}Q_b(\log|c|)$, with $Q_b$ a polynomial of degree two?
The square of the diagonal difference in \cref{eq:roots} already creates
such a candidate, but the other matrix entry and the graph correction
may contribute at the same scale. Compute those contributions before
assigning a coefficient. The binary exact formula
\cref{eq:binary-exact} isolates one part of the task.
\end{question}

\begin{question}[A sharp remainder without a power loss]
Replace $O_\eta(|c|^{3/2-\eta})$ by a bound involving explicit powers of
$\log(1/|c|)$ and, ideally, computable constants. The obstruction is the
use of a global H\"older norm in $PEK$. A localized estimate pairing
$E_c$ with the reduced resolvent may keep track of the small support of
the cusp displacement instead of paying $|c|^{-\alpha}$ twice.
\end{question}

\begin{question}[A complete critical expansion]
Determine whether the pressure and its central spectral data have an
asymptotic expansion in powers of $\sqrt{|c|}$ and $\log|c|$. If so,
give a recursive algorithm and prove closure of the relevant function
class under the reduced-resolvent operations. The present three-term
formula excludes a pure Puiseux series but does not exclude additional
log-periodic or other nonalgebraic contributions at later orders.
\end{question}

\begin{question}[Growing detuning and matched selection]
Extend \cref{eq:detuned-selection} to $|u|\le U(t)$ with $U(t)\to\infty$.
The factors $t^{2ut}$ and the shrinking relative central gap must be
controlled together. Determine a quantitative overlap with fixed
supercritical and subcritical regimes, for both pressure and equilibrium
weights. Taking $u\to\infty$ in the limiting formula alone does not
establish that theorem.
\end{question}

\begin{question}[Two-point regularity and differentiated remainders]
Is the pressure $C^{0,1/2}$ on a full neighborhood of zero? Can its
one-sided derivative, away from exceptional nonzero phases, be expanded
by differentiating \cref{eq:main-pressure}? An asymptotic remainder in
function values does not justify differentiating it. Control between
two moving phases, not just between a phase and zero, is required.
\end{question}

\begin{question}[Quantitative convergence of equilibrium measures]
Obtain explicit rates for the convergence to the selected mixtures,
first in the dual of a fixed H\"older space and then in a transport
metric. Determine whether the atomic weight admits a
$\sqrt{|c|}\log(1/|c|)$ correction and whether its coefficient can be
read from the refined eigenvalues plus localized eigenvector estimates.
This would quantify the finite-phase bias visible in the cosine table.
\end{question}

\begin{question}[The left boundary layer]
The right eigenfunction has the explicit profile
\cref{eq:boundary-detuned}. Determine how the probability $\nu_{s,c}$
concentrates near zero and whether rescaling its restriction yields a
nontrivial measure. The first dual correction $\Xi$ is a finite-part
functional rather than a measure, so it cannot itself be used as a
probability profile. A useful formulation should separate concentrated
mass from the small diffuse tail.
\end{question}

\begin{question}[Other masks and other Jordan multiplicities]
For which nonconsecutive digital masks do a defective isolated
spectral cluster and a finite-part dual functional occur? Establish
conditions under which a logarithmic trace response is universal.
For a Jordan chain of length greater than two, determine the analogous
root scaling, equilibrium simplex, and selection rule. A finite matrix
prediction must be supported by an actual operator realization.
\end{question}

\begin{question}[Certified computation and proof-assistant integration]
Turn the weighted quadrature and graph-compression estimates into
explicit constants, then produce interval enclosures for a nontrivial
phase neighborhood. A proof-assistant route can separate the exact
cotangent and digamma identities, the endpoint integral estimates,
the finite matrix calculation, and the spectral/entropy bridges.
The exact finite-dimensional algebra is a useful first target, but
formalizing only that calculation would not verify the infinite-dimensional
theorem.
\end{question}

\appendix
\section{Proof dependencies and an audit of scope}
\label{app:audit}
The pressure proof follows the dependency chain
\[
\begin{gathered}
\text{telescoping quotient and weighted quadrature}
\Longrightarrow\text{isolated normalized Jordan block}\\
\Longrightarrow\text{graph compression with a quadratic error}\\
\Longrightarrow\text{exact atom/bulk response and finite-part logarithm}\\
\Longrightarrow\text{refined eigenvalues}
\Longrightarrow\text{pressure via positive regularization}.
\end{gathered}
\]
The equilibrium proof additionally uses the first left/right eigenvector
corrections, the cancellation $\Xi(Hf)=\int f$, and the conditional
entropy argument. It does not follow just from the leading pressure
asymptotic or from a pointwise limit of the right eigenfunction.

Three verification boundaries should remain explicit. First, the present
manuscript is not a Lean or Rocq development. Second, collocation output
is not an enclosure of the spectrum of the continuous operator. Third,
a theorem proved in a new research draft is not the same thing as an
independently established priority claim. The source ledger records the
specific repository questions addressed, and the manuscript reconstructs
the analytic ingredients rather than importing unreviewed claims as
unlabeled axioms.

\begin{center}\small
\begin{tabular}{@{}p{.43\textwidth}p{.49\textwidth}@{}}
\toprule
\textbf{Claim} & \textbf{Verification supplied}\\
\midrule
Normalized Jordan chain and gap & Written proof; high-precision checks
of the exact identities.\\[1.5mm]
Logarithmic and linear pressure coefficients & Written finite-part and
matrix proof; independent integral and collocation diagnostics.\\[1.5mm]
Boundary layer and mixture selection & Written norm/duality argument;
left/right collocation checks of a distinguishing observable.\\[1.5mm]
Equilibrium meaning and uniqueness & Written conditional-entropy proof,
including zeros of the weight.\\[1.5mm]
PDF and package integrity & Compilation, reference and font checks,
page rendering, and a build-validation receipt.\\
\bottomrule
\end{tabular}
\end{center}

\clearpage
\begin{thebibliography}{9}
\bibitem{GKS}
P.~Gohlke, M.~Kesseb\"ohmer, and T.~I.~Schindler,
\emph{Generalized Thue--Morse measures: spectral and fractal analysis},
2025, \href{https://arxiv.org/abs/2509.22109v1}{arXiv:2509.22109v1}.

\bibitem{GLS}
P.~Gohlke, G.~Lamprinakis, and J.~Schmeling,
\emph{On a family of singular potentials: Parameter dependence of
thermodynamic characteristics}, 2026,
\href{https://arxiv.org/abs/2603.19001v1}{arXiv:2603.19001v1}.
Theorems 2.2--2.3 and the subsequent regularity question.

\bibitem{Repo}
V.~Reshetnikov, \emph{ProveIt}, repository snapshot
\nolinkurl{1085b506d65e207a05b7e9c861bb1fe88432fe38}, inspected
29 September 2026.
\href{https://github.com/VladimirReshetnikov/ProveIt/tree/1085b506d65e207a05b7e9c861bb1fe88432fe38}{Pinned repository}.

\bibitem{CriticalRepo}
ProveIt research corpus, research draft prepared with ChatGPT for
V.~Reshetnikov,
\emph{Critical Cusps in Digital-Product Pressure: A Jordan collision,
square-root phase scaling, and explicit subcritical response},
29 September 2026. At the snapshot in \cite{Repo}, source
\path{Analysis/FabiusFunction/docs/semi-formalized-research-frontiers/drafts/thue-morse/Thue_Morse_Critical_Pressure/article.tex}.
\href{https://github.com/VladimirReshetnikov/ProveIt/blob/1085b506d65e207a05b7e9c861bb1fe88432fe38/Analysis/FabiusFunction/docs/semi-formalized-research-frontiers/drafts/thue-morse/Thue_Morse_Critical_Pressure/article.tex}{Pinned source}.
The research questions ``The next critical term'' and ``Eigenfunctions,
eigenmeasures, and equilibrium measures'' are the immediate targets.

\bibitem{FractionalRepo}
ProveIt research corpus, research draft prepared with ChatGPT for
V.~Reshetnikov,
\emph{Fractional Cusps at the Atomic Phase: Sharp noninteger regularity
of Thue--Morse pressure and sharp relaxation for digital products in
every base}, 29 September 2026. At the snapshot in \cite{Repo}, source
\path{Analysis/FabiusFunction/docs/semi-formalized-research-frontiers/drafts/thue-morse/Thue_Morse_Fractional_Pressure/article.tex}.
\href{https://github.com/VladimirReshetnikov/ProveIt/blob/1085b506d65e207a05b7e9c861bb1fe88432fe38/Analysis/FabiusFunction/docs/semi-formalized-research-frontiers/drafts/thue-morse/Thue_Morse_Fractional_Pressure/article.tex}{Pinned source}.

\bibitem{SubcriticalRepo}
ProveIt research corpus,
\emph{Thue\_Morse\_Subcritical\_Pressure}, research package,
29 September 2026. At the snapshot in \cite{Repo}, source
\path{Analysis/FabiusFunction/docs/semi-formalized-research-frontiers/drafts/thue-morse/Thue_Morse_Subcritical_Pressure/article.tex}.
\href{https://github.com/VladimirReshetnikov/ProveIt/blob/1085b506d65e207a05b7e9c861bb1fe88432fe38/Analysis/FabiusFunction/docs/semi-formalized-research-frontiers/drafts/thue-morse/Thue_Morse_Subcritical_Pressure/article.tex}{Pinned source}.
Its all-subcritical extension is also recorded in the editorial note in
\cite{CriticalRepo}; no theorem from this package is assumed here.
\end{thebibliography}
\end{document}
